% !TEX root = ../termpaper.tex
% @author Leonhard Lüdeke
%

\section{Methodik und Struktur}

\subsection{Methodisches Vorgehen}

Methodisch kombiniert die Arbeit konzeptionelle Architekturentwicklung mit strukturierter Sicherheitsanalyse.

Zunächst werden auf Basis einschlägiger Literatur sowie der gegebenen Rahmenbedingungen Anforderungen, Randbedingungen und Zielkriterien herausgearbeitet. \cite{georgeff_communication_1983,jennings_commitments_1993,jennings_roadmap_1998,russell_artificial_2022,sycara_distributed_1996,wang_large_2026} Hierbei sind neben den Literatur basierten Rahmenbedingungen, die Unternehmerischen Anforderungen von maßgeblicher Bedeutung. Hierfür werden Fachinterviews mit den zuständigen Abteilungsleitern geführt, um klare Use Cases (Bereichsleitung DSAI= Data Science and Artificial Intelligence) zu definieren und notwendigen Sicherheitsanforderungen (CISO) zu ermitteln. 

Anschließend wird daraus eine Zielarchitektur für ein LLM-basiertes Multi-Agenten-System abgeleitet, in der Rollen, Schnittstellen, Kommunikationsbeziehungen und zentrale Vertrauensgrenzen beschrieben werden.

Im nächsten Schritt wird diese Architektur in Komponenten, Datenflüsse und potenzielle Angriffsflächen zerlegt.

Darauf aufbauend erfolgt das Threat Modelling, das klassische Kategorien verteilter Systeme mit LLM-spezifischen Angriffsmustern verbindet. \cite{wong_adding_2000,tete_threat_2024,krawiecka_extending_2025,noauthor_threat_2026,editor_owasp_nodate}

Die identifizierten Risiken werden anschließend in technische und organisatorische Gegenmaßnahmen übersetzt und hinsichtlich ihrer Wirksamkeit sowie verbleibender Restrisiken diskutiert. \cite{krawiecka_extending_2025,wong_adding_2000,editor_owasp_nodate,shostack2014threat,ferrag_prompt_2026,raza_trism_2026}

Auf diese Weise entsteht ein methodischer Ansatz, der Entwurf, Analyse und Bewertung nicht voneinander trennt, sondern systematisch miteinander verknüpft.

\subsection{Vorläufige Kapitelstruktur}

Die Kapitelstruktur beginnt mit der Darstellung des Problemkontextes über die theoretischen Grundlagen und die Anforderungsanalyse zur Architekturkonzeption, Erstellung einer Sicherheitsanalyse und endet mit einer Evaluation und Beantwortung der Forschungsfragen.

Eine mögliche Struktur umfasst die Kapitel Einleitung, theoretische Grundlagen, methodisches Vorgehen, Anforderungsanalyse, Architekturkonzeption, Threat Modelling, Evaluation, Diskussion und Fazit.

Dabei ist besonders wichtig, dass der Threat-Modelling-Teil nicht als isolierter Sicherheitseinschub erscheint, sondern als integraler Bestandteil der Architekturentwicklung nachvollziehbar wird.